\chapter{Discussion and Conclusions}
\label{cap:discusion}

\section{Discussion of the results}

\subsection{Arbitrage capacity and price tracking (OE2, H1)}

For an isolated constant-product pool, the arbitrage result is a direct consequence of the swap equation: capacity $c$ moves the log price by at most $2\log(1+c)$ per tick until the no-arbitrage band is reached. For the tested log shock of $0.20$, low capacity $c=0.005$ requires 20 ideal ticks to correct the excess outside the band. The simulated reduction in excess area therefore confirms the implementation under interacting order flow, but its direction is not an independent discovery. Its more useful quantitative content is the finite-horizon recovery pattern and the difference between the isolated benchmark and the noisy schedule.

\subsection{Liquidity withdrawal and the H2 sensitivity (OE3)}

The fixed withdrawal rule is activated in every E3 run because each tested shock exceeds its threshold; it is a one-time response and does not represent heterogeneous provider behavior. A replay of the existing seeds clarifies the shallow-cell mechanism. In E1-high static, all 143 observed later departures are single episodes that coincide with the configured panic tick, and the matched no-panic run has no departures. In E3, 98.2--99.6\,\% of recovered paths depart later, usually in multiple episodes, and this pattern remains without panic. A paired control that fixes the arbitrage input cap at its initial-reserve value reduces the shallow E3-minus-static post-recovery excess-area contrast from 0.122--0.124 to 0.0036--0.0097. This indicates that reserve-linked capacity explains much of the contrast in the selected design. The control is post hoc, uses existing seeds, conditions the area comparison on recovery, and covers one depth; it does not isolate a universal causal effect of withdrawal or depth. H2 is therefore a conditional interaction among the threshold rule, panic, order scale, reserves and capacity.

\subsection{Reduced maps and predictive scope (OE4, H3)}

This answers the map clause as limited predictive description, not regime classification: the two-state map improves on persistence in some held-out conditions but does not beat a scalar AR(1), so its second state adds no demonstrated one-step predictive value. The band-aware map's roughly 1.7\,\% mean replay-error improvement is concentrated in one low-capacity/depth condition and depends on the realized future CEX path. It is a conditional replay result, not a real-time forecast.

\section{Conclusions by objective}

\begin{description}
  \item[OE1.] The simulator passed its implementation invariants and its recorded outputs are reproducible. This establishes consistency with the encoded rules.
  \item[OE2.] The ideal capacity correction is derived exactly. Simulation shows the finite-horizon consequences of adding the specified noise, trend and event schedule; it does not provide an independent test that arbitrage moves price toward the reference.
  \item[OE3.] The panic ablation distinguishes single panic-associated E1 departures from recurrent E3 departures that persist without panic. At the tested shallow depth, a fixed-initial-reserve arbitrage cap removes most of the post-recovery E3 contrast. These results identify a mechanism in this simulator, conditional on the tested rule, order schedule and recovery selection; they do not establish a general withdrawal effect.
  \item[OE4.] The two-state map does not improve on scalar AR(1). The band-aware map has a small, localized replay advantage using future realized reference input; evidence for operational prediction is absent.
\end{description}

\section{General conclusion}

Within the specified grid, the isolated-pool calculation gives the ideal correction per event, while the simulated post-perturbation trajectory also depends on reserve depth, order size, agent rules and event order. The revised question asks about these finite-horizon error trajectories rather than market regimes. On that target, the two-state map does not beat scalar AR(1), and the band-aware map's small replay gain is localized and uses the realized future CEX path. The evidence therefore answers the revised question within this simulator but does not establish a real-market effect, an operational forecast or a regime classifier.

The recovery-persistence result is concentrated in E3. Across its nine equally weighted depth/shock cells, 31.2\,\% of paths that achieved sustained recovery later left the band; at depth factor 0.5 the conditional departure share was 98.8\,\%, whereas at factors 1 and 2 it was zero. The 31.2\,\% figure is a design-weighted simulator summary, not an estimated market frequency. In the shallow cells, the additional replay indicates two distinct patterns: E1's later exits are single panic-associated episodes, whereas E3 exits recur even when panic is disabled. A fixed-capacity control substantially attenuates the E3 contrast. Future event replay should record first band entry, two-tick sustained return, later departures, panic/order events, active reserves and available arbitrage capacity separately; these are candidate observables, not validated live monitoring requirements.

The revised H2 estimand excludes the imposed shock and the two recovery-confirmation ticks by measuring excess-error area after sustained return. At depth 0.5, E3 minus E1-high is 0.1229, 0.1239 and 0.1215 for shocks $-0.05$, $-0.10$ and $-0.20$; at depths 1 and 2 the paired contrasts are zero to numerical precision. The departure replay shows why a continuous-flow-only ratio was insufficient: in shallow static E1, $R_{\mathrm{flow}}=0.670<1$, but the one-off panic ratio is $R_{\mathrm{panic}}=1.327>1$. For a conservative one-tick configured pressure measure, use $R_{\max}^{\mathrm{adv}}=1.586$, which includes the maximum aligned noise-plus-panic input; the component ratios distinguish persistent flow from a single impulse. E1's observed departures coincide with this one-off event. In E3, $R_{\mathrm{flow}}=1.208$ after the nominal withdrawal and repeated departures persist without panic. These ratios use configured, aligned orders and approximate reserves; they are explanatory diagnostics, not a validated threshold rule. The shallow fixed-cap control reduces the post-recovery area contrast to 0.0036--0.0097, showing that the original contrast is substantially linked to the reserve-based capacity rule in this design.

For engineering use, these findings justify candidate telemetry fields for later replay experiments: synchronized reference and pool prices, active reserves, signed order flow, available arbitrage capacity, transaction/event order, first band entry, sustained return and subsequent departures. This is a proposal for what to measure in the next study, not evidence that a live monitoring rule is already validated.

The thesis's contribution is methodological and bounded: it derives the constant-product capacity benchmark, tests how selected rules interact in a reproducible simulator, and reports sensitivity and negative model-comparison results. A real-system claim would require synchronized replay of market events with pool fees, gas, latency, finite capital, transaction ordering and competition among searchers. The current manuscript contains no such empirical replay.

\section{Limitations and threats to validity}

The simulator uses an exogenous CEX price, zero-fee swaps, no gas or execution delay, ample agent inventories, one representative agent per role, a constant-product pool and event ticks without calibrated wall-clock duration. It does not represent concentrated-liquidity positions, competition among searchers or MEV. The withdrawal and panic rules are fixed heuristics; the H2 trigger fires in every tested run, and withdrawal changes both depth and the arbitrage cap. The sensitivity and held-out model comparisons remain within this one synthetic rule family. The hybrid comparison additionally uses the realized future CEX path. No calibration, synchronized event replay or external validation is included.
